% Generated draft: requires mathematical and source review.
\documentclass[11pt]{article}
\usepackage[margin=25mm]{geometry}
\usepackage{fontspec}
\setmainfont{Latin Modern Roman}
\usepackage{amsmath,amssymb,mathtools,unicode-math}
\setmathfont{Latin Modern Math}
\usepackage{enumitem,needspace,xurl,hyperref}
\setlength{\parindent}{0pt}
\setlength{\parskip}{5pt}
\setlength{\emergencystretch}{3em}
\newcommand{\N}{\mathbb{N}}
\newcommand{\Z}{\mathbb{Z}}
\newcommand{\Q}{\mathbb{Q}}
\newcommand{\R}{\mathbb{R}}
\newcommand{\C}{\mathbb{C}}
\newcommand{\F}{\mathbb{F}}
\newcommand{\A}{\mathbb{A}}
\newcommand{\PP}{\mathbb{P}}
\newcommand{\E}{\mathbb{E}}
\newcommand{\eps}{\varepsilon}
\newcommand{\dd}{\,\mathrm{d}}
\begin{document}
\section*{UnsolvedMath ID: 30000880}
\textbf{Source problem code:} catalogue entry
\subsection{Definitions and mathematical statement}
The extraction combines three distinct questions.
\begin{enumerate}
\item Let \(G\) be a group, \(K\) a field, and \(t:G\to\operatorname{Aut}(K)\) a group homomorphism. The skew group ring \(A=K^tG\) consists of finite sums \(\sum_g a_g g\), with multiplication
\[
(a g)(b h)=a\,t_g(b)\,(gh).
\]
Does \(A\) have invariant basis number, namely
\[
A^m\cong A^n\text{ as left }A\text{-modules}
\quad\Longrightarrow\quad m=n
\qquad(m,n\in\Z_{\geq0})?
\]
No finiteness or growth hypothesis on \(G\) is imposed in this general question.
\item The source refers to a recursive construction of torsion units \(w_i\) from a torsion element \(a\) and \(b\in G\). It asks what can be said about \(S(a,b)=\{w_i\}\) and the subgroup \(\langle S(a,b)\rangle\) of the relevant unit group. The recurrence, initial term, index range, coefficient ring, and hypotheses on \(a,b\) are absent. Thus the units and subgroup cannot be reconstructed from this extraction, and no replacement recurrence or structural conclusion is asserted.
\item Let \(p>2\) be prime, let \(G\) be a finite nonabelian \(p\)-group, and let \(K\) have characteristic \(p\). Does the ordinary group algebra \(KG\) have no filtered multiplicative basis? With the customary radical-filtration convention, such a basis is a \(K\)-basis \(\mathcal B\) for which products of basis elements lie in \(\mathcal B\cup\{0\}\), and \(\mathcal B\cap J^r\) is a basis of \(J^r\) for every \(r\geq1\), where \(J=\operatorname{rad}(KG)\). This last question concerns the untwisted algebra and all fields of the specified characteristic, not necessarily finite fields.
\end{enumerate}
The radical-filtration convention in the third question should be confirmed in the report. The three questions are retained separately and are not given one uniform unresolved status.
\subsection{Short English statement}
Does every skew group ring of a field have invariant basis number? What is the structure of an unspecified recursively generated set of torsion units? And do nonabelian finite \(p\)-group algebras in odd characteristic lack filtered multiplicative bases?
\subsection{Sources}
\begin{itemize}
\item Local UnsolvedMath registry, permanent ID \(30000880\), source code catalogue entry, ``Basis Numbers and Torsion Units in Skew Group Algebras.'' No individual proposer is supplied.
\item \emph{Mini-Workshop: Arithmetik von Gruppenringen} (2008), Oberwolfach Reports, DOI 10.4171/owr/2007/55, pp. 3209--3240: \url{https://doi.org/10.4171/owr/2007/55}.
\item Supplied reference, ``filtered multiplicative bases'': \url{https://arxiv.org/abs/math/0009008}.
\item Supplied reference, ``further filtered-basis results'': \url{https://arxiv.org/abs/math/0102101}.
\item Supplied literature triage, checked 2026-08-21: status ``partially solved.'' It reports invariant basis number for broad cases including subexponential group growth, and classifications for many finite \(p\)-groups including the negative \(p>2\) filtered-basis case, while describing the combined questions as not fully resolved. This dated assessment derives from supplied OWR curation; time-sensitive status and cited claims require independent verification.
\end{itemize}
\end{document}
